% Options for packages loaded elsewhere
\PassOptionsToPackage{unicode}{hyperref}
\PassOptionsToPackage{hyphens}{url}
%
\documentclass[
  11pt,
]{article}
\usepackage{amsmath,amssymb}
\usepackage{iftex}
\ifPDFTeX
  \usepackage[T1]{fontenc}
  \usepackage[utf8]{inputenc}
  \usepackage{textcomp} % provide euro and other symbols
\else % if luatex or xetex
  \usepackage{unicode-math} % this also loads fontspec
  \defaultfontfeatures{Scale=MatchLowercase}
  \defaultfontfeatures[\rmfamily]{Ligatures=TeX,Scale=1}
\fi
\usepackage{lmodern}
\ifPDFTeX\else
  % xetex/luatex font selection
\fi
% Use upquote if available, for straight quotes in verbatim environments
\IfFileExists{upquote.sty}{\usepackage{upquote}}{}
\IfFileExists{microtype.sty}{% use microtype if available
  \usepackage[]{microtype}
  \UseMicrotypeSet[protrusion]{basicmath} % disable protrusion for tt fonts
}{}
\makeatletter
\@ifundefined{KOMAClassName}{% if non-KOMA class
  \IfFileExists{parskip.sty}{%
    \usepackage{parskip}
  }{% else
    \setlength{\parindent}{0pt}
    \setlength{\parskip}{6pt plus 2pt minus 1pt}}
}{% if KOMA class
  \KOMAoptions{parskip=half}}
\makeatother
\usepackage{xcolor}
\usepackage{graphicx}
\makeatletter
\def\maxwidth{\ifdim\Gin@nat@width>\linewidth\linewidth\else\Gin@nat@width\fi}
\def\maxheight{\ifdim\Gin@nat@height>\textheight\textheight\else\Gin@nat@height\fi}
\makeatother
% Scale images if necessary, so that they will not overflow the page
% margins by default, and it is still possible to overwrite the defaults
% using explicit options in \includegraphics[width, height, ...]{}
\setkeys{Gin}{width=\maxwidth,height=\maxheight,keepaspectratio}
% Set default figure placement to htbp
\makeatletter
\def\fps@figure{htbp}
\makeatother
\setlength{\emergencystretch}{3em} % prevent overfull lines
\providecommand{\tightlist}{%
  \setlength{\itemsep}{0pt}\setlength{\parskip}{0pt}}
\setcounter{secnumdepth}{-\maxdimen} % remove section numbering
\ifLuaTeX
  \usepackage{selnolig}  % disable illegal ligatures
\fi
\usepackage{bookmark}
\IfFileExists{xurl.sty}{\usepackage{xurl}}{} % add URL line breaks if available
\urlstyle{same}
\hypersetup{
  hidelinks,
  pdfcreator={LaTeX via pandoc}}

\author{}
\date{}

\begin{document}

\section{Numerical Experiments}\label{numerical-experiments}

\[
\begin{equation}
\begin{align}
E(y):=
\underbrace{\frac{-1}{2}\int_\Omega |\nabla^2 v|^2 - \nu [v, v]}_{\text{membrane}} +
\underbrace{\frac{1}{2}c_\nu\int_\Omega |\Delta w|^2 - (1-\nu) [w, w]}_{\text{bending}} \nonumber \\  +
\underbrace{\frac{1}{2 }\int_\Omega J^T \nabla^2 v J: \nabla w\otimes \nabla w}_{\text{coupling}} - \gamma \beta^4\underbrace{\int_\Omega p w }_{\text{transv. pressure}} + \beta^2\underbrace{ \sum_i\int_\Omega s_i \delta(x_i) v}_{\text{disclinations}}
\end{align}
\end{equation}
\]

\subsection{1. Numerical Setup}\label{numerical-setup}

\subsubsection{1.1 Variational Problem}\label{variational-problem}

\begin{itemize}
\tightlist
\item
  Brief description of the problem being solved.
\item
  Reference to the finite element discretisation and solver
  configuration.
\end{itemize}

\subsubsection{1.2 Implementation Details}\label{implementation-details}

The numerical solution of the variational problem presented in Section 5
is carried out leveraging the dolfinx/FEniCSx finite element platform,
relying on the PETSc library for efficient linear and nonlinear algebra.
Our implementation, https://github.com/kumiori/disclinations features
high-order polynomial Lagrange finite elements whose order can be
changed at runtime, for flexible accuracy in the discretisation of the
variational principle. More specifically, we use the Newton nonlinear
solver provided by PETSc SNES and without step adjustment (`basic
lineasearch', cf.~manual). The iterative subproblems are solved by
Krylov subspace method (preonly) with LU factorization as preconditioner
(pc\_type: lu).

The numeircal • Stopping tolerance: snes\_stol = 1e-6 • Absolute
tolerance: snes\_atol = 1e-6 • Relative tolerance: snes\_rtol = 1e-6 •
Maximum iterations: snes\_max\_it = 200

Geometry and Mesh

All presented numerical experiments are performed on a circular domain
\(\Omega := B(0, 1)\) (the unit open ball) centered in the origin,
meshed with triangular (Delaunay) elements of typical size
\(\eta = 0.05\). This geometry forms the basis for all parametric
verification computations and numerical experiments. The finite element
discretisation employs third-order Lagrange elements for both the
transverse displacement (\(w\)) and the Airy stress potential (\(v\)),
ensuring sufficient smoothness, allowing for the computation of the
Gaussian curvature, whilst capturing the nonlinear coupling between
in-plane and out-of-plane elastic behaviour.

The numerical stabilisation coefficient \(\alpha\) is fixed to \ldots{}

1,595 nodes, xxxx elements, matrix sizes \textasciitilde{} \ldots K
entries

The non-dimensional formulation involves two tunable parameters The
aspect ratio (\(\beta = R/h \gg 1\)), measuring the plate's slenderness
and the relative transverse pressure (\(p_0/E\)), ratio of applied
pressure to material stiffness. Finally, the disclination strength,
characterised by the Frank angle for kinematically incompatible loads,
is fixed to \(\pm 1\) in the following experiments, to qualitatively
distinguish positive from negative induced membrane stresses,
respectively corresponding to compression and tension.

All computations refer to the non-dimensional formulation.

Nontrivial kinematically incompatible loads, verified against
closed-form analytical solutions.

Mixed loading conditions combining surface pressure and disclinations.

\begin{center}\rule{0.5\linewidth}{0.5pt}\end{center}

\subsection{2. Verification Tests}\label{verification-tests}

Validation experiments provide substantiation that the computational
model accurately represents the mathematical model of the mechanical
system within specified limits of accuracy. These experiments are
functional in establishing the reliability of the numerical framework
and its ability to predict key physical behaviors. The aim of validation
is twofold:

\begin{enumerate}
\def\labelenumi{\arabic{enumi}.}
\tightlist
\item
  Quantitative Accuracy: Demonstrating that the numerical results match
  known solutions within acceptable bounds.
\item
  Qualitative Insights: Ensuring that the model captures the essential
  physical mechanisms and trends inherent to the system.
\end{enumerate}

Validation experiments focus on two configurations: kinematically
compatible loads through surface pressure distributed over the plate and
kinematically incompatible loads with a dipole of disclinations,
symmetrically distributed around the centre of the plate.

\(\rho(x) = \sqrt{x_1^2 + x_2^2}\)

\subsubsection{2.1 Plate with Kinematically Compatible Loads: transverse
pressure}\label{plate-with-kinematically-compatible-loads-transverse-pressure}

rescaling nondimensionalisation, all test collapse to \ldots{}

Let \(\theta = 0\) and
\(p(x)  = p_0\frac{40}{3} (1 - \rho^2(x) )^4 + \frac{16}{3} (11 + \rho^2(x) )\),
then \$w(x) \$ = \(w_0 (1 - \rho^2(x) )^2\) and
\(v(x) =v_0  (-\frac{1}{12} (1 - \rho^2(x) )^2 -\frac{1}{18} (1 - \rho^2(x) )^3 -\frac{1}{24} (1 - \rho^2(x) )^4)\)
solves eq \ldots.

Here, \(p_0=\sqrt{\frac{2D^3}{E h}}\),
\(w_0\sqrt{\frac{2D}{E h}}, v_0 =\frac{D}{h}\) In Figures \ref{} we
compare the approximate solutions of the three finite element
formulations introduced in section \ref{} to the analytical solution
presented above.

Due to the rescaling, the nondimensional formulation involving only
transverse load is independent of the aspect ratio,

the dependence on the aspect ratio is explicit in the rescaling factors
for the fields and the energy

figure \ldots{} we compare the profile of the finite element approximate
fields with the closed form analytic solution, along the segment
\(x_2 = 0\)

Owing to the interpretation of \(v\) as membrane stress potential, its
sign locally informs of the positive vs.~negative character of membrane
strains.

Transverse pressure induces bending which in turn, induces tensile
(positive) membrane strains.

Finally, figure \ldots{} shows the Gaussian curvature, computed
postprocessing the equilibrium solution, as the determinant of the
Hessian of the transverse displacement. Note the inversion of curvature
due to the clamping boundary condition. In the image, the red locus
corresponds to \(x: \kappa(x)=0\)

\begin{itemize}
\tightlist
\item
  \textbf{Setup}: Plate loaded with surface pressure.
\item
  \textbf{Objectives}: Compare numerical results against known analytic
  solutions.
\item
  \textbf{Results}: Error analysis, visualisations.
\end{itemize}

\includegraphics{experiment_parametric_transverse/output/d8cb39/1742dd/transverse.png}
\includegraphics{experiment_parametric_transverse/output/d8cb39/1742dd/transverse-profiles.png}
\includegraphics{experiment_parametric_transverse/output/909867_energy_terms.png}
Caption: remove figure, energy (bending, membrane, and coupling) wrt
closed form computation
\includegraphics{experiment_parametric_transverse/output/909867_error_terms.png}
\includegraphics{experiment_parametric_transverse/output/909867/a945c5/gaussian_curvature.png}
\includegraphics{experiment_parametric_transverse/output/d8cb39/1742dd/gaussian_curvature.png}
Caption: Gaussian curvature, locus of inversion of the curvature Extra:
compare errors varying penalisation parameter?

\subsubsection{2.2 Plate with Kinematically Incompatible Loads, Dipole,
Parametric varying aspect
ratio}\label{plate-with-kinematically-incompatible-loads-dipole-parametric-varying-aspect-ratio}

To verify our numerical implementation for a nontrivial incompatible
loading against closed form analytic solutions, the second setup
involves a plate loaded with a dipole of disclinations symmetrically
arranged around the centre of the plate.

Let \(\theta(x) = \delta(x-0.2 e_1) - \delta(x+0.2 e_1), p(x)= 0\). Then
\(w(x) \equiv 0\) and
\(v(x) = v_0( \rho^2(x) \ln(\rho^2(x)) - \rho^2(x) + 1)\) solves
\eqref{2406171748}. here \(v_0=\frac{E h s}{16\pi}\) We solve
{[}\ldots{]} with \(\gamma = 0\), for varying \(\beta\) between \(10\)
and \(X00\).

\(w=0\) is an admissible solution for the dipole problem

Figure, \ldots, show energy scaling as a function of the aspect ratio,
fourth power of \(\beta\). This behaviour corresponds to the expected
energy scaling, by remarking that the multiplicative factor which scales
the membrane equations - linear when \(w=0\) - is \(\beta^2\). The
stress solution is hence quadratic in \(\beta^2\) thus exponent of the
energy. On the other hand, the bending energy and the coupling energy
are zero.

Figure \ldots{} compares the profiles of approximate finite element
solutions with respect to the closed form analytic solutions.

Figure \ldots{} shows the fields \(w, v\) on the two-dimensional
computational domain, the dots mark the locations of the disclinations
(colour coding),

Figure shows the scalar stress potential associated to the positive
(colour) and negative (colour) disclination, respectively inducing a
locally tensile and compressive state. RK: colours are inverted, FIX

\begin{itemize}
\tightlist
\item
  \textbf{Results}: Error analysis, visualisations.
\end{itemize}

\includegraphics{experiment_parametric_dipole/output/4a96c1_energy_terms.png}
\includegraphics{experiment_parametric_dipole/output/4a96c1/a33f9515b198a07caa4d27cc0b7b367c/dipole.png}
\includegraphics{experiment_parametric_dipole/output/4a96c1/a33f9515b198a07caa4d27cc0b7b367c/dipole-profiles.png}

\begin{center}\rule{0.5\linewidth}{0.5pt}\end{center}

\subsection{3. Parametric Experiments}\label{parametric-experiments}

\subsubsection{3.1 Mixed Loading: Disclination monopole and surface
pressure, Varying
Thickness}\label{mixed-loading-disclination-monopole-and-surface-pressure-varying-thickness}

We perform a series of parametric experiments investigating mixed-mode
loading combining surface pressure and a disclination monopole, as the
plate's aspect ratio is varied whilst keeping the prefactor of the
transverse external load \(\gamma \beta^4\) constant and equal to one.
This implies that, as \(\beta\) varies between \(10\) and \(X00\), the
physical dimensional surface pressure decreases with the inverse fourth
power of the dimensional thickness.

Nondimensional energy plots show the energy scaling of the different
components of the energy. Similarly to the dipole setting, the membrane
energy scales with the fourth power of the aspect ratio, whereas the
non-dimensional bending energy decreases with a (negative) exponent
smaller than \(-2\) {[}show the bound{]} (more precisely, \(-1.65\)).

For large values of the aspect ratio, small numerical oscillations
appear on the bending term, whose order of magnitude is 8 orders of
magnitude smaller than the dominant membrane energy term.

In the dimensionless formulation presented {[}X{]}, \(\beta\) acts as a
gain for the distribution of disclinations, while \(\gamma\) modulates
the external load.

In this section we analyse the mechanics of the system by varying first
\(\beta\) in some appropriate range while keeping \(\gamma\) constant;
and then varying \(\gamma\) with constant \(\beta\). In order to do
this, we provide an estimate for a realistic range of values for both
the dimensionless parameters.

Non-dimensional transverse displacement goes to zero.

We tune \(\overline{p_0}\) to (a multiplier for) the weight of the
plate, i.e.~considering a density \(\rho_{\text{D}} = O(10^3 kg/m^3)\)
(corresponding to a reference density for steel/metal) and denoting by
\(g\) the gravitational acceleration we obtain

\(\overline{p_0} \sim \rho_{\text{D}} g  h \sim 10^2 \text{Pa}\)

\begin{itemize}
\tightlist
\item
  \textbf{Results}: Parametric plots (e.g., energy terms, deformation
  profiles).
\end{itemize}

\includegraphics{experiment_Parametric_Relative_Work_constant_transverse/output/0869de_energy_terms.png}
\includegraphics{experiment_Parametric_Relative_Work_constant_transverse/output/0869de/042eed1759fac538ebfe539ab7e3f60f/state.png}
\includegraphics{experiment_Parametric_Relative_Work_constant_transverse/output/0869de/042eed1759fac538ebfe539ab7e3f60f/state-profiles.png}
\includegraphics{experiment_Parametric_Relative_Work_constant_transverse/output/0869de/a51acb047d059a0c1a4309343a3d6984/state-profiles.png}
\includegraphics{experiment_Parametric_Relative_Work_constant_transverse/output/0869de/bb2d43230517c6468001976846f268b7/state-profiles.png}

For large values of \(a\), \(h\) is small and \(p_0\) growths with the
fourth power of \(h\). For \(p_0 \equiv 0\), \(w \equiv 0\) in
\(B_1(0)\) solves the out-of-plane equation. This is the reason why in
Figure \ref{fig:24108179} the transverse displacement approaches zero as
\(a\) increases. As a consequence, \(v\) in equation \eqref{2410161145}
is only activated by \(\delta(0)\) and hence asymptotically approaches
\(v^0(x)\), solution to the problem

\[
\begin{equation}
\left\{
    \begin{aligned}
        &\Delta^2 v^0 = XXX \delta(0) && \text{ in } B_1(0) \\
        &v^0 = \partial_n v^0 = 0 && \text{ on } \partial B_1(0)
    \end{aligned}
\right.
\end{equation}
\]

\[
\begin{equation}
    v^0(x) = E h s \mathcal{G}_{\Omega}(x) \quad \Rightarrow |\!| v^0 |\!|_{L^2(\Omega)} = \frac{E h s}{30\pi} \approx 10^5 N/m
\end{equation}
\]

where \(\mathcal{G}_{\Omega}\) is the Green function of the biharmonic
problem in \(H^2_0(\Omega)\) and is equal to

\[
\begin{equation}
    \mathcal{G}_{\Omega}(x) = \frac{1}{16\pi} \Big( 1- |x|^2 + |x|^2 \ln (|x|^2) \Big), \qquad x\in \Omega
\end{equation}
\]

\subsubsection{3.2 Mixed Loading: Varying
Pressure}\label{mixed-loading-varying-pressure}

parametric experiment, varying surface pressure in a range that spans
0.1 to 100 times the plate's weight, assuming \ldots{} density

relative interplay of bending vs.~membrane deformation, changing
external loading conditions whilst keeping the aspect ratio and the
power of the disclination constant.

interpreted as .. multiplier of plate's weight

Energy plot shows scaling properties of the energy terms

computations for aspect ratio \(\beta =100\)

both the disclination and the surface pressure induce same sign
curvature this setup shows a relatively weak coupling across the entire
range of the loading parameter, the membrane energy is constant and
equals the energy estimate provided by the closed form solution for the
fundamental monopole solution. On the other hand, the bending energy is
quadratic with respect to the load multiplier, as expected by inspecting
the (strong form of the) bending problem whilst considering the stress
potential a given. As the value of the surface pressure varies, the
stress potential shows a self-similar profile (cf images), whereas the
transverse displacement progressively localises the deformation around
the (clamped) boundary.

\begin{itemize}
\tightlist
\item
  \textbf{Setup}: Surface pressure combined with disclination monopole;
  pressure is varied to span a range \ldots{}
\item
  \textbf{Results}: Parametric plots (e.g., energy terms, deformation
  profiles).
\end{itemize}

\begin{figure}
\centering
\includegraphics{experiment_Parametric_Relative_work_constant_thickness/output/5f1f1d_energy_terms.png}
\caption{image}
\end{figure}

\includegraphics{experiment_02/output/5509b4/mesh.png} typical mesh
(with one indicated disclination at the centre)

\includegraphics{experiment_Parametric_Relative_work_constant_thickness/output/5f1f1d/7e70a6881a231348bb4816b62d437858/state.png}
\includegraphics{experiment_Parametric_Relative_work_constant_thickness/output/5f1f1d/7e70a6881a231348bb4816b62d437858/state-profiles.png}

\begin{center}\rule{0.5\linewidth}{0.5pt}\end{center}

\subsection{4. Complex Experiment: Flower
Configuration}\label{complex-experiment-flower-configuration}

\subsubsection{4.1 Setup}\label{setup}

\begin{itemize}
\tightlist
\item
  \textbf{Configuration}: A flower-like arrangement of disclinations
  mimicking engineered graphene sheets.
\item
  \textbf{Objectives}: Demonstrate the capability of the model for
  complex and nonstandard configurations.
\end{itemize}

\subsubsection{4.2 Results}\label{results}

\begin{itemize}
\tightlist
\item
  Energy analysis, deformation profiles.
\item
  Discussion of insights from this experiment.
\end{itemize}

\begin{center}\rule{0.5\linewidth}{0.5pt}\end{center}

\subsection{5. Summary of Numerical
Results}\label{summary-of-numerical-results}

\begin{itemize}
\tightlist
\item
  Recap key findings from the verification and parametric experiments.
\item
  Discuss implications for the mechanics of thin plates under complex
  loading scenarios.
\end{itemize}

\begin{center}\rule{0.5\linewidth}{0.5pt}\end{center}

\subsubsection{Details}\label{details}

\begin{itemize}
\item
  Software environment: dolfinx/FEniCSx, PETSc.
\item
  Solver configurations (e.g., SNES settings, tolerance). solvers:
  nonlinear: prefix: fvk snes: snes\_type: newtonls snes\_stol: 1e-6
  snes\_atol: 1e-6 snes\_rtol: 1e-6 snes\_max\_it: 200 \#
  snes\_divergence\_tolerance: -1.0 snes\_monitor: '\,' ksp: ksp\_type:
  preonly pc\_type: lu
\item
  Meshes geom\_type: `circle' radius: 1. mesh\_size: 0.05
\end{itemize}

\begin{enumerate}
\def\labelenumi{\arabic{enumi}.}
\tightlist
\item
  Variational Problem: • What is the mathematical formulation? governing
  equations, detaild in sec.~5 boundary conditions, H20
\end{enumerate}

• What are the primary unknowns (e.g., displacements, rotations)?
transverse displacement w, airy membrane stress potential v • What are
the physical parameters involved? Youngs modulus E, thickness h
non-dimensional formulation highlights the role of two non-dimensional
parameters: aspect ratio, beta, R/h relative transverse pressure, p0/E
disclination strength, frank angle (taken either +1 or -1) 2. Finite
Element Model: • What is the type of elements used (e.g., Argyris,
Lagrange)? FEM lagrange, order: 3 • What is the resolution of the mesh
(e.g., element size, refinement levels)?

\begin{enumerate}
\def\labelenumi{\arabic{enumi}.}
\setcounter{enumi}{2}
\tightlist
\item
  Solver Details: • Which solver (e.g., PETSc's SNES) and
  preconditioning methods are used? petsc snes, the linear iterative
  solver is based on the krylov subspace method, preconditioner is an lu
  factorisation
\end{enumerate}

• What are the tolerances for convergence (e.g., relative, absolute
residuals)? solvers: nonlinear: prefix: fvk snes: snes\_type: newtonls
snes\_stol: 1e-6 snes\_atol: 1e-6 snes\_rtol: 1e-6 snes\_max\_it: 200 \#
snes\_divergence\_tolerance: -1.0 snes\_monitor: '\,' ksp: ksp\_type:
preonly pc\_type: lu

• Any specific solver settings or configurations? 4. Physical
Parameters: • Geometry (e.g., radius, thickness). • Material properties
(e.g., elastic moduli). • Loading conditions (e.g., surface pressure, ).

\end{document}
